% ============================================================
% results.tex
% ============================================================

This chapter presents the main empirical findings from the analysis of 12,968 securities class actions filed between 1990 and 2024. It proceeds in six sections. Section~\ref{sec:results_overview} first establishes a nonparametric baseline of case duration and competing outcomes. Building on this foundation, Section~\ref{sec:results_pslra} examines the PSLRA’s associations with resolution timing and type through a progressive sequence of models. Section~\ref{sec:results_circuits} then investigates geographic disparities across federal circuits. Section~\ref{sec:results_characteristics} analyzes the influence of key case-level structural covariates, particularly multidistrict litigation (MDL). Section~\ref{sec:results_diagnostics} evaluates model assumptions and proportional-hazards diagnostics, and Section~\ref{sec:results_summary} synthesizes the principal findings.

%% ============================================================
%% SECTION 5.1 — OVERVIEW
%% ============================================================

\section{Overview: Duration and Competing Outcomes}
\label{sec:results_overview}

Securities class actions resolve relatively quickly in aggregate. Approximately 28\% of cases terminate within the first year and 47\% within two years, driven largely by dismissals during the motion-to-dismiss phase. Resolution slows considerably after year four, with fewer than 7\% of cases remaining pending beyond eight years (Figure~\ref{fig:km_overall}).

\begin{figure}[htbp]
    \centering
    \begin{subfigure}[b]{0.49\textwidth}
        \centering
        \includegraphics[width=\textwidth]{fig_km_overall.pdf}
        \caption{Kaplan-Meier Survival Function}
        \label{fig:km_overall}
    \end{subfigure}
    \hfill
    \begin{subfigure}[b]{0.49\textwidth}
        \centering
        \includegraphics[width=\textwidth]{fig_cif_overall.pdf}
        \caption{Aalen-Johansen CIF}
        \label{fig:cif_overall}
    \end{subfigure}
    \vskip 10pt
    \caption[Overall Baseline: Survival and Cumulative Incidence]{Overall duration and competing-outcome patterns ($n = 12{,}968$). Shaded bands = 95\% pointwise CIs. Left: Kaplan-Meier overall survival function. Right: Aalen-Johansen cumulative incidence functions (CIFs) for settlement and dismissal (Scheme~A).}
    \label{fig:baseline_plots}
\end{figure}

Dismissal is the dominant resolution pathway. Its cumulative incidence reaches approximately 37\% by two years and 51\% by five years, while settlement accumulates more gradually, reaching 29\% at five years and 39\% at eight years (Figure~\ref{fig:cif_overall}). The two curves do not converge, highlighting a structural asymmetry: dismissals concentrate early through pleading-stage motions, whereas settlements occur more slowly among surviving cases.

A naive Kaplan-Meier analysis that censors dismissals would overstate long-run settlement probabilities by treating dismissed cases as though they remained at risk of eventual settlement. The Aalen-Johansen estimator avoids this bias and supplies the appropriate nonparametric baseline against which the PSLRA analyses in the following section are evaluated.

%% ============================================================
%% SECTION 5.2 — THE PSLRA EFFECT
%% ============================================================

\section{The PSLRA and Litigation Outcomes}
\label{sec:results_pslra}

Building on the overall baseline, this section examines how the PSLRA is associated with changes in resolution timing and type. Across multiple specifications, post-PSLRA filing is consistently linked to a higher dismissal hazard, while the association with settlement is directionally negative but considerably less robust. The evidence is developed in five progressive layers: nonparametric patterns, associational estimates, time-varying dynamics, composition-adjusted inference, and robustness checks.

\subsection{Nonparametric Evidence}

The PSLRA coincides with opposing shifts in the two competing processes (Figure~\ref{fig:cif_pslra}). Post-PSLRA, dismissal incidence rises sharply in the first year and reaches 52.1\% by five years (versus 37.8\% pre-PSLRA). Settlement incidence is suppressed, reaching only 26.8\% at five years (versus 51.0\% pre-PSLRA). Gray’s test rejects equality of the CIFs for both outcomes ($p < 0.001$).
\begin{figure}[htbp]
    \centering
    \includegraphics[width=\linewidth]{fig_cif_pslra.pdf}
    \vskip 5pt
    \caption[Cumulative Incidence of Outcomes by PSLRA Regime]{Aalen-Johansen CIFs by PSLRA regime (Pre: $n = 1{,}032$; Post: $n = 11{,}936$). Shaded bands = 95\% pointwise CIs.}
    \label{fig:cif_pslra}
\end{figure}

Table~\ref{tab:cif_horizons} quantifies these differences at key horizons. The divergence is especially pronounced early: one-year dismissal incidence is 25.4\% post-PSLRA versus 15.4\% pre-PSLRA, consistent with the PSLRA's heightened pleading requirements that concentrate
dismissals at the pleading stage. By year five, the ordering of outcomes reverses. Pre-PSLRA, settlement dominated (51.0\% vs.\ 37.8\% dismissal). Post-PSLRA, dismissal (52.1\%) substantially exceeds settlement (26.8\%), creating a 24.2 percentage-point gap.

\begin{table}[htbp]
\centering
\caption{Cumulative Incidence at Key Horizons by PSLRA Regime (\%)}
\label{tab:cif_horizons}
\renewcommand{\arraystretch}{1.2}
\begin{tabular*}{\textwidth}{@{\extracolsep{\fill}} l l r r r r @{}}
\toprule
\textbf{Regime} & \textbf{Outcome} & \textbf{1 yr} & \textbf{2 yr} & \textbf{3 yr} & \textbf{5 yr} \\
\midrule
Pre-PSLRA  & Settlement &  6.4 & 22.4 & 33.3 & 51.0 \\
           & Dismissal  & 15.4 & 25.6 & 29.7 & 37.8 \\
\midrule
Post-PSLRA & Settlement &  3.0 &  9.0 & 16.2 & 26.8 \\
           & Dismissal  & 25.4 & 37.5 & 45.5 & 52.1 \\
\bottomrule
\end{tabular*}
\end{table}

These descriptive patterns are formalized in cause-specific Cox models, which quantify the associations while properly accounting for right-censoring.

\subsection{Associational Evidence: Cause-Specific Cox Models}

Baseline cause-specific Cox models confirm the nonparametric pattern (Table~\ref{tab:cox_baseline}). Post-PSLRA filing is associated with a 44\% lower settlement hazard and a 41\% higher dismissal hazard. Both associations are statistically significant ($p < 0.001$).

\begin{table}[htb]
\centering
\caption{Baseline Cause-Specific Cox Models (PSLRA Only)}
\label{tab:cox_baseline}
\renewcommand{\arraystretch}{1.2}
\begin{tabular*}{\textwidth}{@{\extracolsep{\fill}} l r c r @{}}
\toprule
\textbf{Outcome} & \textbf{HR} & \textbf{95\% CI} & \textbf{$p$-value} \\
\midrule
Settlement & 0.563 & (0.511, 0.621) & $<$0.001 \\
Dismissal  & 1.409 & (1.269, 1.565) & $<$0.001 \\
\bottomrule
\end{tabular*}
\end{table}

The Grambsch-Therneau test rejects proportional hazards for both outcomes ($p < 0.001$), indicating that the PSLRA’s associations vary over time. This motivates the piecewise analysis below; the constant hazard ratios should therefore be interpreted as time-averaged effects of a more complex temporal relationship.

\subsection{Time-Varying Associations}
\label{sec:piecewise}

Piecewise Cox models reveal distinctly outcome-specific temporal patterns (Table~\ref{tab:piecewise}). For dismissal, the PSLRA association is strongest in the first year (HR = 1.789, $p < 0.001$) and declines monotonically, becoming indistinguishable from null beyond two years. For settlement, suppression persists across all intervals, though it is strongest between years one and two (HR = 0.387, $p < 0.001$).

\begin{table}[htbp]
\centering
\caption{Piecewise PSLRA Associations by Outcome and Time Period}
\label{tab:piecewise}
\renewcommand{\arraystretch}{1.2}
\begin{tabular*}{\textwidth}{@{\extracolsep{\fill}} l l r c r @{}}
\toprule
\textbf{Outcome} & \textbf{Period} & \textbf{HR} & \textbf{95\% CI} & \textbf{$p$-value} \\
\midrule
Dismissal  & 0--1 years & 1.789 & (1.520, 2.105) & $<$0.001 \\
           & 1--2 years & 1.255 & (1.016, 1.550) & 0.035    \\
           & 2$+$ years & 1.063 & (0.886, 1.276) & 0.510    \\
\addlinespace[0.4em]
Settlement & 0--1 years & 0.518 & (0.395, 0.678) & $<$0.001 \\
           & 1--2 years & 0.387 & (0.323, 0.464) & $<$0.001 \\
           & 2$+$ years & 0.673 & (0.591, 0.767) & $<$0.001 \\
\bottomrule
\end{tabular*}
\begin{minipage}{\textwidth}
\vskip 4pt
\footnotesize \textit{Note:} Hazard ratios are interval-specific following time-splitting.
\end{minipage}
\end{table}

\subsection{Composition-Adjusted Evidence: IPTW Triangulation}
\label{sec:results_iptw}

Having established these associational patterns, a natural question is how much of the observed divergence survives adjustment for changes in observable case composition. If the
post-PSLRA period sees disproportionately more cases filed
in high-dismissal circuits, for instance, part of the raw
PSLRA association could be an artifact of compositional
shifts rather than a stable period difference under a common
case mix. To address this, we apply inverse probability of treatment weighting (IPTW) using the four-strategy triangulation framework described in Chapter~\ref{ch:methodology}. 

The preferred marginal structural model (MSM) estimate indicates that post-PSLRA filing remains associated with a 26\% lower settlement hazard (HR = 0.741) and a 52\% higher dismissal hazard (HR = 1.519), even after reweighting the pre-PSLRA cohort to match the post-PSLRA covariate distribution (Table~\ref{tab:iptw_triangulation}).

\begin{table}[htbp]
\centering
\caption[PSLRA Hazard Ratios: Four Estimation Strategies]{PSLRA Hazard Ratios Under Four Estimation Strategies ($n = 12{,}866$)}
\label{tab:iptw_triangulation}
\renewcommand{\arraystretch}{1.2}
\begin{tabular*}{\textwidth}{@{\extracolsep{\fill}} l r c r c @{}}
\toprule
& \multicolumn{2}{c}{\textbf{Settlement}} & \multicolumn{2}{c}{\textbf{Dismissal}} \\
\cmidrule(lr){2-3} \cmidrule(lr){4-5}
\textbf{Strategy} & \textbf{HR} & \textbf{95\% CI} & \textbf{HR} & \textbf{95\% CI} \\
\midrule
1.\ Unadjusted                    & 0.557 & (0.505, 0.615) & 1.414 & (1.273, 1.570) \\
2.\ Regression-Adjusted           & 0.784 & (0.709, 0.867) & 1.661 & (1.493, 1.847) \\
3.\ IPTW + Regression Covariates  & 0.780 & (0.682, 0.892) & 1.755 & (1.549, 1.987) \\
4.\ Marginal Structural Model     & 0.741 & (0.632, 0.868) & 1.519 & (1.335, 1.729) \\
\bottomrule
\end{tabular*}
\begin{minipage}{\textwidth}
\vskip 4pt
\footnotesize \textit{Note:} All estimates $p < 0.001$. IPTW uses ATT weights with 99th-percentile trimming. Rows~3 and~4 rely on robust (Lin-Wei sandwich) standard errors. These intervals quantify sampling uncertainty under the weighted specification but do not resolve concerns regarding overlap or post-treatment bias in the weighting variables. Row~4 (MSM) is the preferred composition-adjusted specification as it avoids outcome-model assumptions; comparisons across rows are point-estimate contrasts rather than formal decompositions and should be read as sensitivity results.
\end{minipage}
\end{table}

The IPTW-weighted CIFs (Figure~\ref{fig:cif_weighted}) show that the settlement gap narrows after reweighting, while the dismissal gap persists. These estimates adjust for observable compositional shifts but cannot rule out unmeasured temporal confounding or post-treatment bias in covariates such as circuit or MDL status.

\begin{figure}[htbp]
    \centering
    \begin{subfigure}[b]{0.49\textwidth}
        \centering
        \includegraphics[width=\textwidth]{fig_cif_weighted_settlement.pdf}
        \caption{Settlement}
        \label{fig:cif_weighted_settlement}
    \end{subfigure}
    \hfill
    \begin{subfigure}[b]{0.49\textwidth}
        \centering
        \includegraphics[width=\textwidth]{fig_cif_weighted_dismissal.pdf}
        \caption{Dismissal}
        \label{fig:cif_weighted_dismissal}
    \end{subfigure}
    \vskip 10pt
    \caption[IPTW-Weighted Cumulative Incidence by PSLRA Regime]{IPTW-weighted Aalen-Johansen CIFs by PSLRA regime (ATT weights, 99th-percentile trimming). The pre-PSLRA cohort is reweighted to match the post-PSLRA observable case mix.}
    \label{fig:cif_weighted}
\end{figure}


\subsection{Robustness Checks}
\label{sec:results_robustness}

The directional findings hold across seven alternative specifications (Figure~\ref{fig:robustness}, Table~\ref{tab:robustness}). Five observations warrant particular emphasis.

\begin{figure}[htbp]
    \centering
    \includegraphics[width=\linewidth]{fig_robustness_hr.pdf}
    \vskip 5pt
    \caption[PSLRA Hazard Ratios Across Robustness Specifications]{PSLRA hazard ratios across robustness specifications. Each row is a separate Cox specification. Note the log scale on the horizontal axis.}
    \label{fig:robustness}
\end{figure}

\begin{table}[htb]
\centering
\caption{PSLRA Hazard Ratios Across Robustness Specifications}
\label{tab:robustness}
\renewcommand{\arraystretch}{1.2}
\begin{tabular*}{\textwidth}{@{\extracolsep{\fill}} l l r c r @{}}
\toprule
\textbf{Specification} & \textbf{Outcome} & \textbf{HR} & \textbf{95\% CI} & \textbf{$p$-value} \\
\midrule
Scheme A: Narrow        & Settlement & 0.563 & (0.511, 0.621) & $<$0.001 \\
                        & Dismissal  & 1.409 & (1.269, 1.565) & $<$0.001 \\
\addlinespace[0.5em]
Scheme B: Intermediate  & Settlement & 0.799 & (0.730, 0.874) & $<$0.001 \\
                        & Dismissal  & 1.183 & (1.053, 1.330) & 0.005 \\
\addlinespace[0.5em]
Scheme C: Broad          & Settlement & 0.788 & (0.722, 0.859) & $<$0.001 \\
                        & Dismissal  & 1.183 & (1.053, 1.330) & 0.005 \\
\addlinespace[0.5em]
Excl.\ post-2020        & Settlement & 0.558 & (0.506, 0.616) & $<$0.001 \\
                        & Dismissal  & 1.333 & (1.200, 1.481) & $<$0.001 \\
\addlinespace[0.5em]
Second Circuit only     & Settlement & 1.052 & (0.757, 1.462) & 0.762    \\
                        & Dismissal  & 1.941 & (1.387, 2.715) & $<$0.001 \\
\addlinespace[0.5em]
Ninth Circuit only      & Settlement & 0.957 & (0.803, 1.140) & 0.622    \\
                        & Dismissal  & 1.575 & (1.329, 1.865) & $<$0.001 \\
\addlinespace[0.5em]
Time-trend (spline)     & Settlement & 0.714 & (0.548, 0.932) & 0.013    \\
                        & Dismissal  & 3.515 & (2.730, 4.525) & $<$0.001 \\
\bottomrule
\end{tabular*}
\end{table}

First, alternative disposition coding schemes (Schemes B and C) that reclassify Code~12 (voluntary dismissal) and Code~5 (consent
judgment) as settlements produce more attenuated settlement HRs ($\approx 0.79$, $p < 0.001$) compared with the baseline Scheme A (HR = 0.563, $p < 0.001$). The settlement association remains statistically significant under all schemes, while dismissal HRs moderate from 1.409 to 1.183.

Second, excluding post-2020 cases (which carry high censoring) has virtually no effect on the estimates, indicating that recently filed cases do not distort the results.

Third, circuit-specific sub-models provide the most important qualification of the settlement finding. Within the Second Circuit the PSLRA settlement HR crosses unity (HR = 1.052, $p = 0.762$), and within the Ninth Circuit the association is non-significant (HR = 0.957, $p = 0.622$). By contrast, the dismissal association remains strong and significant in both circuits (Second: HR = 1.941; Ninth: HR = 1.575). This pattern suggests that part of the pooled settlement association reflects between-circuit compositional differences rather than a uniform within-circuit reform effect.

Fourth, the most parametrically flexible check—a natural cubic spline time trend (df = 3)—produces the largest dismissal HR in the robustness set (HR = 3.515, $p < 0.001$), while the settlement HR attenuates to 0.714 ($p = 0.013$). The PSLRA indicator exhibits a variance inflation factor (VIF) of 5.32 with the spline basis, indicating substantial but not perfect overlap with secular trends. This specification should be interpreted as an upper bound on the dismissal association rather than a central estimate---it shows that the positive association survives flexible time-trend control, not that it is free of calendar-time confounding. The attenuation of the settlement HR from 0.563 suggests that part of the raw settlement suppression is similarly confounded with secular trends.

Fifth, across all seven specifications the dismissal HR ranges from 1.183 to 3.515 (always positive and significant), while the settlement HR ranges from 0.558 to 1.052. The dismissal direction is unanimous; the settlement direction is consistent in six of seven models but fragile under circuit-specific and cluster-robust scrutiny.

To further probe pre-existing trends, we implement two temporal identification checks. A clean-window analysis restricted to the narrow 1993--1998 period surrounding the PSLRA’s effective date ($n = 1{,}235$) yields a dismissal HR of 1.482 ($p < 0.001$) and a settlement HR of 0.724 ($p < 0.001$), both directionally consistent with the full-sample results. 

As a falsification check, we construct a placebo ``law'' at January~1, 1992 and re-estimate on the pre-PSLRA sample only ($n = 1{,}032$, spanning 1990--1995). The placebo produces a dismissal HR below unity (HR = 0.733, $p = 0.006$) and a settlement HR above unity (HR = 1.638, $p < 0.001$), indicating a downward pre-existing trend in dismissal rates during the early 1990s---the opposite direction from
the PSLRA's subsequent upward association. That directional
contrast rules out a simple monotone story in which an
upward pre-trend mechanically inflates the later dismissal
estimate, but it also shows that period-specific dynamics
were already moving before the law changed. Together, these checks reinforce the decision to frame the estimates as robust associations rather than strictly causal effects.

Circuit-specific results within these robustness checks, together with the flexible spline time-trend specification, point to the broader importance of judicial geography. 

%% ============================================================
%% SECTION 5.3 — GEOGRAPHIC DISPARITY
%% ============================================================

\section{Geographic Disparity Across Federal Circuits}
\label{sec:results_circuits}

Circuit geography emerges as one of the strongest predictors of securities class action outcomes. Settlement hazard ratios relative to the Second Circuit range from 1.1 to 4.3 times across circuits, indicating that judicial venue substantially shapes both the likelihood and timing of resolution.

\subsection{Cumulative Incidence by Circuit}

Cumulative incidence functions reveal substantial heterogeneity that differs markedly by outcome (Figure~\ref{fig:cif_circuit}).

\begin{figure}[htbp]
    \centering
    \begin{subfigure}[b]{0.49\textwidth}
        \centering
        \includegraphics[width=\textwidth]{fig_cif_circuit_dismissal.pdf}
        \caption{Dismissal Incidence by Circuit}
        \label{fig:cif_circuit_dismissal}
    \end{subfigure}
    \hfill
    \begin{subfigure}[b]{0.49\textwidth}
        \centering
        \includegraphics[width=\textwidth]{fig_cif_circuit_settlement.pdf}
        \caption{Settlement Incidence by Circuit}
        \label{fig:cif_circuit_settlement}
    \end{subfigure}
    \vskip 10pt
    \caption[Cumulative Incidence of Outcomes by Judicial Circuit]{Aalen-Johansen CIFs by circuit (top 4 by volume plus Sixth Circuit). Gray's test rejects equality for both outcomes ($p < 0.001$).}
    \label{fig:cif_circuit}
\end{figure}

The Fifth Circuit has the highest cumulative dismissal incidence among the plotted circuits, reaching 34.7\% by year one and 56.0\% by year six---marginally above the Second Circuit's trajectory. The Third Circuit is the lowest-dismissal curve at 42.4\% by year six.

Settlement incidence reverses this ordering. The Third Circuit leads, reaching 29.4\% by year three, while the Second Circuit is the lowest among the plotted circuits at 22.3\% by year six. The two figures together characterize the Second Circuit as a comparatively low-settlement, high-dismissal venue.

\subsection{Circuit Effects in Cause-Specific Cox Models}


Every circuit has a settlement hazard 1.1 to 4.3 times that of the Second Circuit 
(Table~\ref{tab:cox_circuit}), with the Third Circuit showing the largest gap 
(HR~$= 4.326$). The Second Circuit is the lowest-settlement 
venue in the data, though the magnitude of this gap narrows once the PSLRA is 
controlled for. On the dismissal side, the Fourth Circuit stands out: its dismissal 
hazard is only 37\% of the Second Circuit's (HR~$= 0.367$), 
the lowest of any circuit; the Tenth Circuit also falls below the reference 
(HR~$= 0.745$). Circuit geography alone achieves moderate discrimination 
(C-index: settlement $= 0.646$, dismissal $= 0.569$).

\begin{table}[htbp]
\centering
\caption[Circuit Effects: Cause-Specific Cox Models]{Circuit Effects:
Cause-Specific Cox Models}
\label{tab:cox_circuit}
\renewcommand{\arraystretch}{1.2}
\begin{tabular*}{\textwidth}{@{\extracolsep{\fill}} l r c r c @{}}
\toprule
& \multicolumn{2}{c}{\textbf{Settlement}} & \multicolumn{2}{c}{\textbf{Dismissal}} \\
\cmidrule(lr){2-3} \cmidrule(lr){4-5}
\textbf{Covariate} & \textbf{HR} & \textbf{95\% CI} & \textbf{HR} & \textbf{95\% CI} \\
\midrule
Post-PSLRA    & 0.754 & (0.683, 0.834) & 1.558 & (1.401, 1.733) \\
First Circ.   & 2.248 & (1.751, 2.886) & 1.442 & (1.239, 1.680) \\
Third Circ.   & 4.326 & (3.768, 4.965) & 1.043 & (0.918, 1.186) \\
Fourth Circ.  & 1.134 & (1.008, 1.275) & 0.367 & (0.318, 0.422) \\
Fifth Circ.   & 1.986 & (1.664, 2.371) & 1.396 & (1.247, 1.564) \\
Sixth Circ.   & 3.830 & (3.216, 4.562) & 0.995 & (0.841, 1.178) \\
Seventh Circ. & 2.536 & (2.140, 3.005) & 1.500 & (1.323, 1.701) \\
Eighth Circ.  & 3.672 & (3.070, 4.393) & 1.051 & (0.903, 1.223) \\
Ninth Circ.   & 2.451 & (2.233, 2.691) & 1.183 & (1.108, 1.262) \\
Tenth Circ.   & 3.233 & (2.785, 3.751) & 0.745 & (0.637, 0.872) \\
Eleventh Circ.& 2.165 & (1.803, 2.600) & 1.528 & (1.355, 1.723) \\
\bottomrule
\end{tabular*}
\begin{minipage}{\textwidth}
\vskip 4pt
\footnotesize \textit{Note:} Reference category: Second Circuit
($n = 12{,}952$). DC and Federal Circuits omitted ($n < 50$ each).
\end{minipage}
\end{table}

\subsection{PSLRA \texorpdfstring{$\times$}{x} Circuit Interaction}

Likelihood ratio tests comparing the additive circuit-PSLRA specification
against a full interaction model reject a uniform PSLRA association for
both outcomes: settlement ($\chi^2 = 51.45$, df~$= 5$,
$p = 7.0 \times 10^{-10}$) and dismissal ($\chi^2 = 71.87$, df~$= 5$,
$p = 4.2 \times 10^{-14}$). Table~\ref{tab:interaction} presents results
restricted to the five largest circuits plus the Second Circuit reference
($n = 10{,}652$). The main PSLRA coefficient captures the reform's
association within the Second Circuit; each interaction term gives the
multiplicative modification for the named circuit.\footnote{Circuit-specific
PSLRA hazard ratios are obtained by multiplying the main PSLRA HR by the
relevant interaction HR (e.g., Third Circuit settlement:
$1.411 \times 0.377 \approx 0.532$).}

\begin{table}[htbp]
\centering
\caption[PSLRA $\times$ Circuit Interaction Models]{PSLRA $\times$ Circuit Interaction Models}
\label{tab:interaction}
\renewcommand{\arraystretch}{1.2}
\begin{tabular*}{\textwidth}{@{\extracolsep{\fill}} l r c r r c r @{}}
\toprule
& \multicolumn{3}{c}{\textbf{Settlement}} & \multicolumn{3}{c}{\textbf{Dismissal}} \\
\cmidrule(lr){2-4} \cmidrule(lr){5-7}
\textbf{Term} & \textbf{HR} & \textbf{95\% CI} & \textbf{$p$-value}
              & \textbf{HR} & \textbf{95\% CI} & \textbf{$p$-value} \\
\midrule
PSLRA (main)       & 1.411 & (1.016, 1.959)  & 0.040
                   & 2.660 & (1.904, 3.716)  & $< 0.001$ \\
\addlinespace[0.4em]
\textit{Interactions} & & & & & & \\
$\times$ Third    & 0.377 & (0.250, 0.567) & $< 0.001$
                  & 2.325 & (1.214, 4.453) & 0.011 \\
$\times$ Fourth   & 0.745 & (0.308, 1.804) & 0.514
                  & 0.145 & (0.081, 0.259) & $< 0.001$ \\
$\times$ Fifth    & 0.380 & (0.229, 0.629) & $< 0.001$
                  & 0.447 & (0.284, 0.703) & $< 0.001$ \\
$\times$ Ninth    & 0.761 & (0.525, 1.101) & 0.148
                  & 0.589 & (0.405, 0.855) & 0.005 \\
$\times$ Eleventh & 0.235 & (0.137, 0.403) & $< 0.001$
                  & 0.678 & (0.395, 1.165) & 0.160 \\
\midrule
\multicolumn{7}{l}{C-index (in-sample): 0.690 (settlement), 0.607 (dismissal)} \\
\bottomrule
\end{tabular*}
\begin{minipage}{\textwidth}
\vskip 4pt
\footnotesize \textit{Note:} Reference category: Second Circuit
($n = 10{,}652$). In-sample C-indices are for within-specification
comparison only; not directly comparable to held-out metrics in
Tables~\ref{tab:performance_cox} and~\ref{tab:performance_fg}.
\end{minipage}
\end{table}

For settlement, the PSLRA main effect is marginally significant in the
Second Circuit (HR~$= 1.411$, $p = 0.040$), while interaction terms for
the Third, Fifth, and Eleventh Circuits are all well below unity
(HR~$= 0.377$, $0.380$, and $0.235$, respectively; all $p < 0.001$), indicating that
settlement suppression was concentrated in the Second Circuit and
substantially attenuated elsewhere.

The dismissal pattern is stronger and more consistent across circuits. The
main PSLRA term implies a 166\% higher dismissal hazard in the Second
Circuit post-reform (HR~$= 2.660$, $p < 0.001$). Attenuation is most pronounced in the Fourth (HR = 0.145, $p < 0.001$) and Fifth Circuits (HR = 0.447, $p < 0.001$), consistent with prior work documenting cross-circuit variation in PSLRA pleading practice and the Second Circuit's role as a comparatively demanding securities venue \citep{perino2003, choi2007}.

\subsection{Shared Frailty Sensitivity Analysis}
\label{sec:frailty}

Residual unobserved circuit-level heterogeneity persists
after conditioning on observed covariates, though it is
modest in magnitude. We fit shared frailty models
(Gaussian random intercepts on the log-hazard scale) as a sensitivity analysis, noting that the
11 circuits fall below the 20--30 clusters typically
recommended for reliable variance estimation. The extended
settlement frailty variance ($\hat{\theta} = 0.130$)
suggests that fixed-effect circuit indicators capture most
but not all circuit-level variation in settlement timing; the
extended dismissal frailty variance is smaller still
($\hat{\theta} = 0.040$), indicating that the dismissal
hazard is comparatively uniform across circuits after
covariate adjustment (Table~\ref{tab:frailty}).

\begin{table}[htbp]
\centering
\caption{PSLRA Hazard Ratio: Fixed-Effects, Frailty, and Cluster-Robust}
\label{tab:frailty}
\renewcommand{\arraystretch}{1.2}
\begin{tabular*}{\textwidth}{@{\extracolsep{\fill}} l r c r @{}}
\toprule
\textbf{Model} & \textbf{HR} & \textbf{95\% CI} & \textbf{$p$-value} \\
\midrule
\multicolumn{4}{l}{\textit{Settlement}} \\
\quad Extended FE (naive SE) & 0.784 & (0.709, 0.867) & $2.4 \times 10^{-6}$ \\
\quad Extended FE (cluster-robust SE) & 0.784 & (0.559, 1.099) & 0.158 \\
\quad Extended RE (frailty, $\hat{\theta} = 0.130$) & 0.783 & (0.707, 0.866) & $2.0 \times 10^{-6}$ \\
\addlinespace[0.5em]
\multicolumn{4}{l}{\textit{Dismissal}} \\
\quad Extended FE (naive SE) & 1.661 & (1.493, 1.847) & $8.6 \times 10^{-21}$ \\
\quad Extended FE (cluster-robust SE) & 1.661 & (1.298, 2.126) & $5.5 \times 10^{-5}$ \\
\quad Extended RE (frailty, $\hat{\theta} = 0.040$) & 1.660 & (1.493, 1.847) & $9.1 \times 10^{-21}$ \\
\bottomrule
\end{tabular*}
\begin{minipage}{\textwidth}
\vskip 4pt
\footnotesize \textit{Note:} All models include the full extended covariate set ($n = 12{,}866$). FE = fixed effects (circuit dummies). RE = random effects (Gaussian frailty on circuit). Cluster-robust SE clusters on circuit (11 clusters). $\hat{\theta}$ is the estimated variance of the Gaussian random intercept.
\end{minipage}
\end{table}

The PSLRA point estimates are virtually identical across all
three estimation strategies (settlement HR $\approx 0.78$,
dismissal HR $\approx 1.66$), confirming that the
fixed-effects specification adequately adjusts for
between-circuit confounding. However, cluster-robust standard
errors are substantially wider than naive standard errors:
the settlement confidence interval widens to 0.559--1.099,
with $p$ rising from $2.4 \times 10^{-6}$ to 0.158---no
longer statistically significant at conventional levels.
The dismissal interval widens to 1.298--2.126 but remains
highly significant ($p = 5.5 \times 10^{-5}$). The
settlement result's loss of significance under cluster-robust
inference reinforces the fragility documented in the IPTW
and circuit-specific analyses: the PSLRA's settlement
association depends on pooling across circuits and does not
survive conservative variance estimation that accounts for
within-circuit correlation.

\section{Case Characteristics: MDL, Origin, and Statutory Basis}
\label{sec:results_characteristics}

Beyond the PSLRA and circuit geography, case-level characteristics---most
notably MDL consolidation---exert strong and heterogeneous effects on
resolution dynamics. Extended cause-specific Cox models show that MDL
assignment is associated with substantially lower hazards for both outcomes,
while Fine-Gray subdistribution models reveal a settlement-side sign
reversal that cause-specific analysis alone obscures.

\subsection{Extended Cause-Specific Cox Models}

The PSLRA associations are robust to the addition of origin, MDL status,
jurisdictional basis, and statutory basis: the settlement HR is 0.784 and the dismissal HR is 1.661 in the extended specification ($n = 12{,}866$;
Tables~\ref{tab:cox_ext_settle} and~\ref{tab:cox_ext_dismiss}). These
estimates are close to the circuit-only specification (0.754 and 1.558,
respectively), confirming that the PSLRA associations are not artifacts of
compositional differences in origin, MDL status, or statutory basis across
regimes.


\begin{table}[tp]
\centering
\caption[Extended Cause-Specific Cox Model: Settlement]{Extended Cause-Specific Cox Model: Settlement ($n = 12{,}866$)}
\label{tab:cox_ext_settle}
\renewcommand{\arraystretch}{1.2}
\begin{tabular*}{\textwidth}{@{\extracolsep{\fill}} l r c @{}}
\toprule
\textbf{Covariate} & \textbf{HR} & \textbf{95\% CI} \\
\midrule
Post-PSLRA            & 0.784 & (0.709, 0.867) \\
\addlinespace[0.5em]
\textit{Circuit (Ref: Second Circuit)} & & \\
First Circuit         &  1.782 & (1.385, 2.292) \\
Third Circuit         &  3.481 & (3.021, 4.012) \\
Fourth Circuit        &  1.926 & (1.592, 2.331) \\
Fifth Circuit         &  1.722 & (1.439, 2.061) \\
Sixth Circuit         &  3.086 & (2.583, 3.689) \\
Seventh Circuit       &  2.011 & (1.691, 2.392) \\
Eighth Circuit        &  3.062 & (2.554, 3.670) \\
Ninth Circuit         &  1.944 & (1.762, 2.145) \\
Tenth Circuit         &  2.805 & (2.398, 3.281) \\
Eleventh Circuit      &  1.625 & (1.346, 1.961) \\
\addlinespace[0.5em]
\textit{Case Characteristics} & & \\
Origin: MDL           &  0.397 & (0.330, 0.478) \\
Origin: Other         &  1.720 & (1.543, 1.918) \\
MDL Flag              &  0.490 & (0.425, 0.566) \\
Federal Question      &  2.830 & (1.702, 4.706) \\
\addlinespace[0.5em]
\textit{Statutory Basis (Ref: 10(b))} & & \\
Section 11            &  1.091 & (0.979, 1.217) \\
Other/Both Statutory  &  0.994 & (0.902, 1.095) \\
Missing               &  0.831 & (0.600, 1.151) \\
\bottomrule
\end{tabular*}
\end{table}

\begin{table}[tp]
\centering
\caption[Extended Cause-Specific Cox Model: Dismissal]{Extended Cause-Specific Cox Model: Dismissal ($n = 12{,}866$)}
\label{tab:cox_ext_dismiss}
\renewcommand{\arraystretch}{1.2}
\begin{tabular*}{\textwidth}{@{\extracolsep{\fill}} l r c @{}}
\toprule
\textbf{Covariate} & \textbf{HR} & \textbf{95\% CI} \\
\midrule
Post-PSLRA            & 1.661 & (1.493, 1.847) \\
\addlinespace[0.5em]
\textit{Circuit (Ref: Second Circuit)} & & \\
First Circuit         & 1.194 & (1.025, 1.390) \\
Third Circuit         & 0.870 & (0.766, 0.989) \\
Fourth Circuit        & 0.643 & (0.527, 0.783) \\
Fifth Circuit         & 1.202 & (1.073, 1.347) \\
Sixth Circuit         & 0.833 & (0.702, 0.988) \\
Seventh Circuit       & 1.200 & (1.058, 1.361) \\
Eighth Circuit        & 0.910 & (0.782, 1.058) \\
Ninth Circuit         & 0.968 & (0.907, 1.033) \\
Tenth Circuit         & 0.727 & (0.621, 0.851) \\
Eleventh Circuit      & 1.222 & (1.084, 1.378) \\
\addlinespace[0.5em]
\textit{Case Characteristics} & & \\
Origin: MDL           & 0.323 & (0.261, 0.400) \\
Origin: Other         & 0.857 & (0.777, 0.945) \\
MDL Flag              & 0.106 & (0.084, 0.134) \\
Federal Question      & 2.148 & (1.368, 3.373) \\
\addlinespace[0.5em]
\textit{Statutory Basis (Ref: 10(b))} & & \\
Section 11            & 1.167 & (1.074, 1.269) \\
Other/Both Statutory  & 1.097 & (1.019, 1.181) \\
Missing               & 0.851 & (0.646, 1.122) \\
\bottomrule
\end{tabular*}
\end{table}


MDL consolidation is associated with large suppression of both outcome
hazards. Cases assigned to an MDL docket have a 51\% lower settlement
hazard (HR $= 0.490$) and an 89\% lower dismissal
hazard (HR $= 0.106$). The asymmetric
magnitude---dismissal is suppressed far more than settlement---indicates
that MDL consolidation does not merely slow resolution uniformly but
fundamentally shifts the balance between competing outcomes, a pattern
examined further in the Fine-Gray comparison below. Federal question
jurisdiction is associated with higher hazards for both outcomes
($\text{HR}_\text{settle} = 2.830$, $\text{HR}_\text{dismiss} = 2.148$).
Section~11 cases have a modestly higher dismissal hazard than
Section~10(b) cases (HR $= 1.167$).

Grambsch-Therneau global tests reject proportional hazards in both
extended models (settlement: $\chi^2 = 597.4$, dismissal:
$\chi^2 = 361.8$; both $p < 0.001$). The time-averaged interpretation noted
in Section~\ref{sec:results_pslra} applies here as well.


\subsection{Cause-Specific versus Subdistribution Models}

The cause-specific and subdistribution frameworks diverge sharply for
covariates that suppress both competing hazards simultaneously---most
dramatically for MDL consolidation, where the two estimands carry opposite
signs (Table~\ref{tab:fg_comparison}). The cause-specific HR holds the
competing pathway fixed; the Fine-Gray subdistribution hazard ratio (SHR) directly reflects cumulative incidence.

\begin{table}[htbp]
\centering
\caption[Cause-Specific vs.\ Subdistribution Hazard Ratios]{Cause-Specific
  vs.\ Subdistribution Hazard Ratios: Selected Covariates}
\label{tab:fg_comparison}
\renewcommand{\arraystretch}{1.2}
\begin{tabular*}{\textwidth}{@{\extracolsep{\fill}} l l r r c r @{}}
\toprule
\textbf{Outcome} & \textbf{Covariate} & \textbf{CS HR} & \textbf{FG SHR} & \textbf{FG 95\% CI} & \textbf{FG $p$} \\
\midrule
Settlement & Post-PSLRA & 0.784 & 0.503 & (0.452, 0.560) & $< 0.001$ \\
           & MDL Flag   & 0.490 & 2.468 & (2.241, 2.717) & $< 0.001$ \\
\midrule
Dismissal  & Post-PSLRA & 1.661 & 1.719 & (1.549, 1.908) & $< 0.001$ \\
           & MDL Flag   & 0.106 & 0.174 & (0.138, 0.219) & $< 0.001$ \\
\bottomrule
\end{tabular*}
\begin{minipage}{\textwidth}
\vskip 4pt
\footnotesize \textit{Note:} CS = cause-specific hazard ratio; FG =
Fine-Gray subdistribution hazard ratio. FG $p$ uses robust sandwich
standard errors.
\end{minipage}
\end{table}

For the PSLRA, Fine-Gray subdistribution hazard ratios are more extreme
than cause-specific ratios for both outcomes (settlement SHR $= 0.503$
vs.\ CS HR $= 0.784$; dismissal SHR $= 1.719$ vs.\ CS HR $= 1.661$). The
amplification on the settlement side occurs because the PSLRA increases
dismissals, shrinking the pool of cases available to settle; the
subdistribution model captures this indirect pathway while the
cause-specific model holds it fixed.

The MDL comparison is the most analytically revealing divergence. Despite
a substantially lower cause-specific settlement hazard (CS HR $= 0.490$),
the Fine-Gray SHR for settlement is significantly above unity
(SHR $= 2.468$). The sign reversal
is explained by the competing pathway: MDL cases have an extremely
low dismissal hazard (CS HR $= 0.106$), so the pool of MDL cases available
to settle remains large throughout follow-up. A lower instantaneous
settlement rate therefore produces a higher cumulative settlement
probability because so few cases exit via dismissal along the way. 

\section{Model Diagnostics and Validation}
\label{sec:results_diagnostics}

This section evaluates the assumptions underlying the models presented in
Sections~\ref{sec:results_pslra}--\ref{sec:results_characteristics} and
assesses out-of-sample predictive performance.

\subsection{Proportional Hazards Assessment}
\label{subsec:ph_diagnostics}

As noted in Section~\ref{sec:results_pslra}, PH violations do not
invalidate the Cox model but imply that constant hazard ratios are
time-averaged summaries. Table~\ref{tab:ph_tests} reports the
Grambsch--Therneau test applied to the extended cause-specific models.

\begin{table}[htbp]
\centering
\caption{Proportional Hazards Tests: Extended Cause-Specific Cox Models}
\label{tab:ph_tests}
\renewcommand{\arraystretch}{1.2}
\begin{tabular*}{\textwidth}{@{\extracolsep{\fill}} l r r r r @{}}
\toprule
& \multicolumn{2}{c}{\textbf{Settlement}} &
  \multicolumn{2}{c}{\textbf{Dismissal}} \\
\cmidrule(lr){2-3} \cmidrule(lr){4-5}
\textbf{Covariate} & $\chi^2$ & $p$ & $\chi^2$ & $p$ \\
\midrule
\texttt{post\_pslra}   & 92.35  & $< 0.001$ &  14.67 & $< 0.001$ \\
\texttt{circuit\_f}     & 268.9  & $< 0.001$ & 177.1 & $< 0.001$ \\
\texttt{origin\_cat}    & 241.2  & $< 0.001$ & 132.0 & $< 0.001$ \\
\texttt{mdl\_flag}      & 197.3  & $< 0.001$ &  1.053 & 0.305 \\
\texttt{juris\_fq}      &  0.004 & 0.947     &   2.668 & 0.102 \\
\texttt{stat\_basis\_f} &  59.76 & $< 0.001$ &  63.26 & $< 0.001$ \\
\addlinespace[0.3em]
\textbf{GLOBAL}        & 597.4  & $< 0.001$ & 361.8 & $< 0.001$ \\
\bottomrule
\end{tabular*}
\end{table}


Both models reject the global proportional hazards null ($\chi^2 = 597.4$
and $361.8$, respectively; both $p < 0.001$). This is expected in a sample
of $n = 12{,}866$ spanning 34 years: even modest departures from
proportionality are detected with high power. The PSLRA indicator violates
proportional hazards in both models---settlement ($\chi^2 = 92.35$,
$p < 0.001$) and dismissal ($\chi^2 = 14.67$, $p < 0.001$)---confirming
that the PSLRA's association with both outcomes varies over case duration,
consistent with the piecewise analysis in Section~\ref{sec:results_pslra}.
Circuit, origin, and statutory basis covariates also violate
proportionality in both models, reflecting their evolving role over
multi-decade follow-up.

The Fine-Gray models reject proportional hazards similarly. The post-PSLRA
term rejects in both the settlement subdistribution model ($\chi^2 = 54.92$,
$p < 0.001$; global: $\chi^2 = 1463.8$, $p < 0.001$) and
the dismissal model ($\chi^2 = 11.78$, $p < 0.001$; global:
$\chi^2 = 346.5$, $p < 0.001$).

Figure~\ref{fig:schoenfeld} displays scaled Schoenfeld residuals for the
PSLRA indicator. In the settlement model, the smoothed trajectory drifts
upward toward the constant-HR baseline, indicating that the
settlement-suppression association weakens at longer durations. In the
dismissal model, the trajectory trends downward over time, consistent with
the piecewise finding that the PSLRA's dismissal association is strongest
in the first year and decays thereafter.

\begin{figure}[htbp]
    \centering
    \begin{subfigure}[b]{0.48\textwidth}
        \centering
        \includegraphics[width=\textwidth]{fig_schoenfeld_settlement_pslra.pdf}
        \caption{Settlement Model}
        \label{fig:schoenfeld_settlement}
    \end{subfigure}
    \hfill
    \begin{subfigure}[b]{0.48\textwidth}
        \centering
        \includegraphics[width=\textwidth]{fig_schoenfeld_dismissal_pslra.pdf}
        \caption{Dismissal Model}
        \label{fig:schoenfeld_dismissal}
    \end{subfigure}
    \vskip 6pt
    \caption[Scaled Schoenfeld Residuals for the PSLRA Indicator]{Scaled Schoenfeld residuals for the PSLRA
    indicator in the extended settlement (left) and dismissal
    (right) models. The solid curve is the smoothed
    time-varying coefficient $\hat{\beta}(t)$; the dashed
    red line marks the constant fitted coefficient. The visible banding is an expected artifact of a binary PSLRA indicator under highly imbalanced post-to-pre risk sets; the curve tracks the time-varying coefficient across those bands.}
    \label{fig:schoenfeld}
\end{figure}

\subsection{IPTW Covariate Balance}
\label{subsec:iptw_balance}

The validity of the composition-adjusted estimates in
Section~\ref{sec:results_iptw} depends on the IPTW weights achieving
adequate balance between the pre- and post-PSLRA cohorts on all observed
covariates. ATT weights with 99th-percentile trimming (cap $\approx 43.5$)
reweight the pre-PSLRA cohort ($n = 1{,}031$) to match the post-PSLRA
covariate distribution ($n = 11{,}835$).

After weighting, all 19 covariate rows achieve $|\text{SMD}| < 0.10$
(Figure~\ref{fig:iptw_balance}); the adjusted propensity-score distance is
0.053 and the largest residual imbalance is $\text{SMD} = 0.035$
(Section~11 statutory basis). The effective sample size of the reweighted
pre-PSLRA cohort is 578 (56.0\% of the unweighted $n$), indicating
moderate information loss but sufficient power, as confirmed by all eight
IPTW hazard ratios achieving $p < 0.001$ (Table~\ref{tab:iptw_triangulation}).

\begin{figure}[htbp]
\centering
\includegraphics[width=0.75\textwidth]{fig_iptw_balance.pdf}
\caption[IPTW Covariate Balance: Love Plot]{Love plot of standardized mean
differences before and after IPTW weighting (ATT weights, 99th-percentile
trimming).}
\label{fig:iptw_balance}
\end{figure}

Figure~\ref{fig:ps_overlap} illustrates the positivity challenge underlying
this analysis. Overlap is concentrated in the upper range of the
propensity-score distribution (roughly 0.75--0.95), reflecting the 11.5:1
post-to-pre imbalance inherent to the legislative date cutoff; trimming at
the 99th percentile discards the most extreme pre-PSLRA reweightings.

\begin{figure}[!htb]
\centering
\includegraphics[width=0.85\linewidth]{fig_ps_overlap.pdf}
\caption[Propensity Score Overlap by PSLRA Regime]{Mirror density plot of
propensity scores for pre-PSLRA (lower, reflected) and post-PSLRA (upper)
cases by PSLRA regime.}
\label{fig:ps_overlap}
\end{figure}

Two caveats apply. First, balance on observed covariates does not guarantee
balance on unobserved confounders; unmeasured changes in judicial attitudes,
plaintiff-bar sophistication, or economic conditions across the 34-year
window remain unaddressed. Second, several covariates---circuit, MDL
status, case origin---may themselves be consequences of the PSLRA rather
than pre-treatment confounders, in which case adjusting for them conditions
on potential downstream channels of the reform. The composition-adjusted
estimates should therefore be interpreted as the association remaining after
conditioning on those observables, not as a lower bound on the total PSLRA
effect.

\subsection{Predictive Discrimination}
\label{subsec:discrimination}

Tables~\ref{tab:performance_cox} and~\ref{tab:performance_fg} compare
out-of-sample predictive discrimination on the held-out 30\% test set
($n = 3{,}861$) using the concordance index (C-index), time-dependent AUC,
and integrated Brier score (IBS). These metrics apply to the cause-specific
Cox and Fine-Gray models only; the IPTW and frailty models estimate
treatment effects and sensitivity parameters rather than individual-level
predictions.

\begin{table}[htbp]
\centering
\caption[Predictive Discrimination: Cause-Specific Cox Models]{Cause-Specific Cox Models: Held-Out Test Set ($n = 3{,}861$)}
\label{tab:performance_cox}
\renewcommand{\arraystretch}{1.2}
\begin{tabular*}{\textwidth}{@{\extracolsep{\fill}} l r r r r r r @{}}
\toprule
& & \multicolumn{4}{c}{\textbf{AUC}} & \\
\cmidrule(lr){3-6}
\textbf{Model} & \textbf{C-index} & \textbf{1yr} & \textbf{2yr}
  & \textbf{3yr} & \textbf{5yr} & \textbf{IBS} \\
\midrule
Cox (Settlement) & 0.679 & 0.732 & 0.708 & 0.705 & 0.793 & 0.121 \\
Cox (Dismissal)  & 0.597 & 0.595 & 0.618 & 0.643 & 0.775 & 0.184 \\
\bottomrule
\end{tabular*}
\begin{minipage}{\textwidth}
\vskip 4pt
\footnotesize \textit{Note:} IBS computed over $[0, 5]$ years
(competing events treated as censored). Models use a reduced covariate
set excluding \texttt{stat\_basis\_f} due to complete separation.
\end{minipage}
\end{table}

\begin{table}[htbp]
\centering
\caption[Predictive Discrimination: Fine-Gray Models]{Fine-Gray Models: Held-Out Test Set ($n = 3{,}861$)}
\label{tab:performance_fg}
\renewcommand{\arraystretch}{1.2}
\begin{tabular*}{\textwidth}{@{\extracolsep{\fill}} l r r r r r @{}}
\toprule
& & \multicolumn{4}{c}{\textbf{AUC}} \\
\cmidrule(lr){3-6}
\textbf{Model} & \textbf{C-index} & \textbf{1yr} & \textbf{2yr}
  & \textbf{3yr} & \textbf{5yr} \\
\midrule
Fine-Gray (Settlement) & 0.555 & 0.659 & 0.601 & 0.531 & 0.377 \\
Fine-Gray (Dismissal)  & 0.595 & 0.592 & 0.611 & 0.633 & 0.770 \\
\bottomrule
\end{tabular*}
\begin{minipage}{\textwidth}
\vskip 4pt
\footnotesize \textit{Note:} C-index approximated by treating competing
events as censored. IBS not reported (subdistribution predictions
not supported). Time-dependent AUC computed with IPCW weights.
\end{minipage}
\end{table}

The predictive picture is uneven. The cause-specific Cox settlement model
shows the strongest discrimination overall, with a non-monotone AUC
trajectory weakest at three years and strongest at five (C-index = 0.679;
AUC@5yr = 0.793). The dismissal-side Cox model is weaker, suggesting that
unobserved case-merit factors remain important. Fine-Gray performance is
lower overall, and the settlement subdistribution model deteriorates
sharply at longer horizons: its AUC falls from 0.531 at three years to
0.377 at five, crossing below the no-skill threshold between those
horizons. This sub-0.5 AUC reflects a known limitation of the approximation adopted here: as noted in Chapter~\ref{ch:methodology}, the Fine-Gray C-index is computed by treating competing events as censored rather than retaining them in the subdistribution risk set. At short horizons this approximation is relatively benign, but at longer horizons the risk-set distortion compounds, causing the linear predictor to rank surviving cases in the wrong direction relative to the true cumulative incidence. The result should therefore be interpreted as a failure of the approximation at long follow-up rather than as evidence that the underlying subdistribution model is inversely predictive. Cox IBS values of 0.121 (settlement) and 0.184 (dismissal) confirm improvement over naive marginal prediction but indicate substantial residual uncertainty.

\begin{figure}[htbp]
\centering
\includegraphics[width=0.9\linewidth]{fig_auc_over_time.pdf}
\caption[Time-Dependent AUC: Cox and Fine-Gray Models]{Time-dependent AUC
at 1, 2, 3, and 5-year horizons for cause-specific Cox and Fine-Gray
models on the held-out test set. Dotted line marks the no-skill classifier
(AUC $= 0.5$).}
\label{fig:auc_over_time}
\end{figure}

%% ============================================================
%% SECTION 5.6 — SUMMARY
%% ============================================================
\section{Summary of Findings}
\label{sec:results_summary}

Two broad conclusions emerge from the analyses. First, the PSLRA's association with higher dismissal rates is consistent
across specifications, particularly in the early pleading stage. This association persists after composition adjustment and
across a wide range of robustness specifications, including clean-window
and placebo-cutoff checks that address pre-existing trends.

Second, the association with lower settlement rates is directionally
consistent in most models but considerably more fragile. It attenuates
under flexible time-trend controls, loses significance in major circuit
sub-samples and under cluster-robust inference, and proves sensitive to
disposition coding choices.

Beyond the PSLRA, circuit geography and MDL consolidation exert powerful,
heterogeneous effects on litigation outcomes. Competing-risks methods
illuminate important patterns---especially the MDL sign reversal---that
single-outcome survival models obscure.

These findings, together with the identification limits that qualify causal
interpretation, are interpreted in greater depth in
Chapter~\ref{ch:discussion}.
